\documentclass[10pt,twocolumn,letterpaper]{article}

\usepackage[
  letterpaper,
  top=1in,
  bottom=1.125in,
  left=0.8125in,
  right=0.8125in,
  columnsep=0.3125in
]{geometry}

\usepackage{times}
\usepackage{microtype}
\usepackage{booktabs}
\usepackage{amsmath}
\usepackage{url}
\usepackage{hyperref}
\usepackage{enumitem}
\usepackage{fancyhdr}
\usepackage{multicol}
\usepackage{tabularx}

\pagestyle{fancy}
\fancyhf{}
\renewcommand{\headrulewidth}{0pt}
\cfoot{\thepage}
\setlength{\footskip}{0.375in}

\hypersetup{colorlinks=true,linkcolor=black,urlcolor=blue,citecolor=black}

\setlength{\parindent}{0pt}
\setlength{\parskip}{3pt}

% Tighten section spacing
\usepackage{titlesec}
\titlespacing*{\section}{0pt}{6pt}{3pt}
\titlespacing*{\subsection}{0pt}{4pt}{2pt}
\titlespacing*{\subsubsection}{0pt}{3pt}{1pt}

\title{\vspace{-0.6cm}\textbf{Animal Breed Classification: Preliminary Report}}

\author{
  Arvin Effat \quad Hanna Kasi \quad Jue Li \quad Yannick Radulescu \\[1pt]
  Ludwig-Maximilians-Universit\"{a}t M\"{u}nchen \\[1pt]
  \small\texttt{\{Ar.Effat, Ha.Kasi, Jue.Li, Y.Radulescu\}@campus.lmu.de}
}
\date{}

\begin{document}
\maketitle
\thispagestyle{fancy}

\begin{abstract}
\small
This preliminary report covers our literature research, dataset evaluation, and coarse system plan for a two-stage cat and dog breed classifier. The system classifies images into 20 breeds (10 cat, 10 dog) or returns $-1$ for rejected inputs. We use YOLO for detection, a CNN trained from scratch for classification, and Grad-CAM for post-hoc explainability.
\end{abstract}

%------------------------------------------------------------------
\section{Introduction}

Fine-grained breed recognition is challenging: breeds share overall body structure and differ only in subtle coat patterns and facial proportions. Our pipeline first localizes the primary animal with a pretrained YOLO detector, then classifies the cropped region with a CNN trained from random initialization. A confidence threshold gates the output; low-confidence predictions return $-1$.

%------------------------------------------------------------------
\section{Literature Research}

\textbf{YOLO for Animal Phenotyping~\cite{yolo_animal_2025}.}
A 2025 survey confirms YOLO as a well-validated, computationally efficient detector for animal-related CV tasks. We restrict its output to cat/dog classes to reduce false positives and inference overhead.

\textbf{GCViT for Breed Recognition~\cite{gcvit_2026}.}
A pretrained Global Context Vision Transformer (GCViT) classifies breeds from the full image, attending to both local and global context. GCViT-Tiny achieves 92.00\% test / 94.54\% validation accuracy with rotation, flip, and brightness augmentation. Strong results, but no explicit localization of discriminative parts. \textit{Not adopted}: we prefer an explicit detect-then-crop pipeline.

\textbf{Multi-Stage ViT~\cite{transfg_2021}.}
Three-stage approach: (1)~full-image ViT features + attention-based bounding box; (2)~cropped object features; (3)~discriminative part crops as independent inputs. Well-suited to visually similar breeds but too compute-intensive. \textit{Not adopted.}

\textbf{Modified VGG-16 Pipeline~\cite{vgg16_github}.}
Two-stage system (cat-vs-dog classifier $\rightarrow$ breed classifier) closely aligned with our architecture. Used as an implementation reference.

\textbf{Grad-CAM~\cite{gradcam_matlab,gradcam_omnixai}.}
Gradient-weighted Class Activation Mapping visualizes which spatial regions drive network decisions. Applied post-training to verify focus on anatomically relevant features rather than background artifacts (see Section~\ref{sec:xai}).

%------------------------------------------------------------------
\section{Data Research}

\subsection{Target Classes}

Twenty classes are required: cat breeds 0--9 and dog breeds 10--19.

\smallskip
\noindent
\begin{tabular}{@{}cl@{\hspace{10pt}}cl@{}}
\toprule
\textbf{Idx} & \textbf{Cat breed} & \textbf{Idx} & \textbf{Dog breed} \\
\midrule
0 & Abyssinian        & 10 & Beagle \\
1 & Bengal            & 11 & Pug \\
2 & Birman            & 12 & Boxer \\
3 & Bombay            & 13 & Shiba Inu \\
4 & British Shorthair & 14 & Samoyed \\
5 & Maine Coon        & 15 & Golden Retriever \\
6 & Ragdoll           & 16 & German Shepherd \\
7 & Sphynx            & 17 & Siberian Husky \\
8 & Tabby             & 18 & Dalmatian \\
9 & Tiger Cat         & 19 & Rottweiler \\
\bottomrule
\end{tabular}

\subsection{Dataset Evaluation}

Three public datasets were assessed for class coverage, size, balance, quality, and reject-class availability.

\smallskip
\noindent\textbf{[D1] AtharvaTaras/Dog-Breeds-Dataset}~\cite{dataset_dog} --- FCI dog breeds. We extracted the 10 target breeds: 35+ images per class, balanced. Varied backgrounds and poses. No quality issues. Non-target FCI breeds serve as reject-class examples.

\smallskip
\noindent\textbf{[D2] ma7555/cat-breeds-dataset}~\cite{dataset_cat} --- 70,899 images across 67 breeds with per-image metadata (age, gender, coat). Abyssinian, Bengal, Bombay confirmed; Tabby and Tiger Cat require extraction from broader labels (``Domestic Shorthair'', ``Tuxedo''). Heavily unbalanced; contains false labels, duplicates, and watermarks. A cleaning pass is required before use.

\smallskip
\noindent\textbf{[D3] zippyz Oxford-style}~\cite{dataset_oxford} --- 37 classes ($\sim$200 images/class, $\sim$7,400 total) modeled on the Oxford-IIIT Pet benchmark. All 10 dog breeds present; 8/10 cat breeds present (\textbf{Tabby} and \textbf{Tiger Cat} absent). No quality issues. Non-target breeds provide clean reject examples.

\smallskip
\noindent
\begin{tabular}{@{}lcccc@{}}
\toprule
\textbf{Criterion} & \textbf{[D1]} & \textbf{[D2]} & \textbf{[D3]} \\
\midrule
Species         & Dog    & Cat     & Both \\
Images          & 350+   & 70{,}899 & ${\sim}7{,}400$ \\
Balanced        & \checkmark & $\times$ & ${\sim}$ \\
All 10 dogs     & \checkmark & ---     & \checkmark \\
All 10 cats     & ---    & partial & 8/10 \\
Reject data     & \checkmark & \checkmark & \checkmark \\
Quality issues  & None   & Several & None \\
\bottomrule
\end{tabular}

\smallskip
D1 and D3 are the primary training sources. D2 supplements the two missing cat classes after cleaning.

%------------------------------------------------------------------
\section{System Design}

\subsection{Inference Pipeline}

\begin{enumerate}[leftmargin=*, topsep=1pt, itemsep=1pt, parsep=0pt]
  \item \textbf{Resize \& normalize:} First, the raw input image is resized to match the static input dimensions required by the object detection network. Pixel values are normalized to a consistent range (typically $[0, 1]$ or standardized using dataset-specific mean and standard deviation vectors) to ensure numerical stability.
  \item \textbf{YOLO detection:} The preprocessed image is passed through a YOLO (You Only Look Once) detector. This step is configured to suppress all non-animal predictions. If the detector fails to localize either a cat or a dog with sufficient confidence, the pipeline triggers an early exit and immediately returns $-1$, filtering out irrelevant background images early on.
  \item \textbf{Primary subject selection}: In scenarios where multiple animals or bounding boxes are detected within a single frame, a heuristic selection mechanism is applied. The system retains only the largest bounding box by area, assuming it represents the primary subject of interest.
  \item \textbf{Cropping and Background Discarding:} The image is cropped tightly around the coordinates of the selected bounding box. All peripheral image data outside this region is discarded to eliminate background noise that could distract the subsequent classifier
  \item[\textit{opt.}] \textit{SAM segmentation}: To further refine the input, an optional Segment-Anything Model (SAM) can be integrated. SAM generates a pixel-precise foreground mask of the animal. While this effectively isolates the subject from complex textures, it introduces significant computational overhead and latency.
  \item \textbf{CNN classification}: The cropped (and potentially masked) image is fed into the fine-grained Convolutional Neural Network (CNN). The network processes the specialized features of the animal and outputs a probability distribution across the target space $\in\{0,\ldots,19\}$, alongside a corresponding confidence score.
  \item \textbf{Threshold}: confidence below threshold $\Rightarrow$ return $-1$.
\end{enumerate}

\subsection{Architecture \& Training}

\textbf{Detector:} utilizes a pretrained YOLO framework where the output layers are filtered so that only the 'cat' and 'dog' classes from the original dataset can trigger a detection.

\textbf{Classifier:} is a separate CNN architecture. Diverging from standard transfer learning approaches, this network is trained from random initialization using only the tightly cropped bounding boxes from the training set. Depending on the remaining hardware and compute budget, an optional spatial or channel-wise attention module may be integrated into the CNN to help the model focus on highly localized distinct features (e.g., ear shapes or eye distances).

\textbf{Optimizer:} SGD with momentum, weight decay, and a learning-rate scheduler. Cross-entropy loss; gradients via backpropagation.

\textbf{Augmentation:} To artificially expand the dataset and improve generalization, several data augmentation techniques are applied during the training phase. These include geometric transformations such as horizontal flips and random rotations, as well as photometric adjustments like brightness variations. Crucially, color jittering and grayscale transformations are strictly avoided or heavily constrained. Preserving intact color channels (RGB) is mandatory, as the visual identity of several target breeds depends heavily on distinct coat colors and patterns (e.g., the solid black of a Bombay cat or the characteristic spots of a Dalmatian).

\textbf{Data split:} The dataset provided by the SEP organizers is strictly partitioned to ensure unbiased evaluation. The full training split is utilized exclusively for model optimization and weight adjustments. A separate validation set is used during development for hyperparameter tuning and threshold selection. The final performance metrics of the system are determined using a completely held-out test set, ensuring the model's ability to generalize to unseen real-world data.

%------------------------------------------------------------------
\section{Explainability}
\label{sec:xai}

After training, Grad-CAM is applied to selected validation images. The procedure: (1)~preprocess and forward-pass the image; (2)~compute gradients of the target class score w.r.t.\ the last convolutional layer's feature maps; (3)~derive per-map importance weights via global average pooling; (4)~form a weighted sum of feature maps with ReLU activation; (5)~upsample and overlay as a heat map. Visualizations are reported for both correct predictions and failure cases to characterize model behavior qualitatively.

%------------------------------------------------------------------
\begin{thebibliography}{9}
\small
\bibitem{yolo_animal_2025}
Anon. \textit{A Review of YOLO Algorithms in Animal Phenotyping Applications}. 2025.

\bibitem{gcvit_2026}
Anon. \textit{Fine-Grained Cat Breed Recognition with GCViT}. arXiv:2602.07534, 2026.

\bibitem{transfg_2021}
Anon. \textit{Multi-Stage Vision Transformer for Fine-Grained Recognition}. arXiv:2106.10587, 2021.

\bibitem{vgg16_github}
xigaoli. \textit{Modified VGG-16 for Animal Breed Classification}. \url{github.com/xigaoli/animal-breed-classification}, 2025.

\bibitem{gradcam_matlab}
MathWorks. \textit{Grad-CAM Reveals the Why Behind Deep Learning Decisions}. MATLAB \& Simulink, 2024.

\bibitem{gradcam_omnixai}
Salesforce. \textit{Grad-CAM for Image Classification (PyTorch)}. OmniXAI Docs, 2025.

\bibitem{dataset_oxford}
zippyz. \textit{Cats and Dogs Breeds Classification Oxford Dataset}. Kaggle, 2025.

\bibitem{dataset_cat}
ma7555. \textit{Cat Breeds Dataset}. Kaggle, 2025.

\bibitem{dataset_dog}
AtharvaTaras. \textit{Dog Breeds Dataset (FCI)}. GitHub, 2025.
\end{thebibliography}

\end{document}